% Replay: prompt-cache invalidation in agentic coding workloads.
%
% Every figure in this document is traceable to a dated file under
% docs/evidence/. Where a figure comes from a different corpus than the one
% before it, the corpus is named in the sentence that carries it. That is not
% padding: an earlier draft of this paper merged four corpora into one and
% reported a 98.8% -> 4.2% "correction" between two quantities that do not
% share a denominator. See Section 6.
%
% Compile: pdflatex replay-preprint.tex  (twice, for the table references)

\documentclass[11pt,twocolumn,letterpaper]{article}
\usepackage[utf8]{inputenc}
\usepackage{geometry}
\geometry{top=1in,bottom=1in,left=0.75in,right=0.75in}
\usepackage{amsmath,amssymb,amsfonts}
\usepackage{booktabs}
\usepackage{graphicx}
\usepackage{url}
\usepackage{microtype}
\usepackage{abstract}

\title{\textbf{Prompt Cache Invalidation in Agentic Coding Workloads:\\
An Attribution Study on Unsolicited Transcripts}}

\author{
  Daniel Saito \\
  RedRobot KK, Tokyo \\
  \texttt{https://github.com/RedRobotKK/Replay}
}
\date{\small{11 September 2026}}

\begin{document}

\twocolumn[
  \begin{abstract}
  Prompt caching is the dominant cost lever in long-running agentic coding
  sessions, and its failures are silent: when a cached prefix is invalidated the
  provider re-bills the historical window at write prices, returns a normal
  response, and raises nothing. We report an attribution study over local agent
  transcripts that were not written for measurement. On a corpus of 1{,}506
  transcripts we isolate 735 cache breaks accounting for 31.26M re-billed
  tokens, and attribute each to one of five causes. On a separate corpus of 60
  main-lane transcripts comprising 26{,}675 consecutive assistant-turn pairs, we
  test the keepalive hypothesis --- that periodic pings preserve cache warmth ---
  and find it addresses the wrong failure: 87\% of cache-creation tokens are
  emitted after gaps shorter than five minutes, during which no published
  time-to-live had elapsed. We further report a measurement artifact in which
  session-wide prefix comparison misattributed cause in fan-out workloads, and
  document the correction. We argue that turn-level differential attribution,
  with explicit provenance tiers, is a prerequisite for cost claims about
  agentic systems, and that aggregate consumption tracking cannot substitute
  for it.
  \vspace{0.5cm}
  \end{abstract}
]

\section{Introduction}

Stateful language-model pricing discounts the re-reading of a cached context
prefix, typically by an order of magnitude relative to writing it. In
autonomous coding agents the context history grows with every turn, so the
retention of that prefix, rather than the length of any single prompt,
determines the economics of a session.

Cache invalidation in these systems is opaque by construction. A provider that
cannot serve a cached prefix does not fail; it re-computes the prefix, bills it
at write rates, and returns a well-formed response. The event is observable
only in the usage counters attached to that response, and only by a reader who
is comparing them against the preceding request. In practice the first signal an
engineer receives is an invoice.

This paper describes an attribution method implemented in an open-source
diagnostic engine, \texttt{Replay}, and reports what it finds on transcripts
produced by ordinary work. We distinguish throughout between aggregate
consumption --- how many tokens a session billed --- and differential
attribution --- which turn re-billed, and against which predecessor.

\subsection{Contributions}

\begin{enumerate}
  \item A five-way cause taxonomy for cache invalidation, with a measured
    distribution over 735 breaks (Section 3).
  \item A falsification of the keepalive mitigation for this workload class,
    with the counterfactual saving priced (Section 4).
  \item A documented measurement artifact in fan-out attribution, and the
    isolation requirement it implies (Section 5).
  \item A provenance discipline that makes the preceding three auditable, and
    a report of what it cost us to violate it (Section 6).
\end{enumerate}

\section{Method}

\subsection{Corpora}

Four corpora appear in this paper. They are distinct, and we name the relevant
one at each figure rather than reporting a single aggregate, because the
quantities they support are not interchangeable.

\begin{table}[h]
\centering
\caption{Corpora. One transcript is one agent lane; a session that spawns
sub-agents contributes several.}
\label{tab:corpora}
\vskip 0.1in
\small
\begin{tabular}{llr}
\toprule
\textbf{Label} & \textbf{Extent} & \textbf{Date} \\
\midrule
$C_{\text{cause}}$ & 1{,}506 transcripts & 2026-09-06 \\
$C_{\text{fid}}$   & 1{,}751 transcripts, 116 sessions & 2026-09-10 \\
$C_{\text{gap}}$   & 60 transcripts, 26{,}675 pairs & 2026-09-10 \\
$C_{\text{lane}}$  & 60 requests, 5 sessions, 16 lanes & 2026-09-06 \\
\bottomrule
\end{tabular}
\end{table}

None of these transcripts was generated for this study. They are the residue of
unrelated engineering work on a single machine, which bounds external validity
(Section 7) and removes one familiar bias: no workload was shaped to the
hypothesis, because the hypothesis postdates the workload.

\subsection{Attribution}

For each consecutive pair of requests within one lane, the engine compares the
serialised prefix, the declared tool set, the system block, the model
identifier, and the inter-request interval. A break is recorded when the
response reports cache-creation tokens that a cache hit would have made
unnecessary. Cause is assigned by the first discriminating comparison, and a
pair for which no comparison discriminates is recorded as \emph{unknown} rather
than assigned to a default.

\subsection{Provenance tiers}

Every figure the engine emits is labelled with one of three tiers, and the
labels are enforced in the reporting path rather than applied editorially:

\begin{itemize}
  \item \textbf{Measured.} Usage counters returned by the provider for a
    specific request.
  \item \textbf{Estimated.} Derived from a byte-to-token fit, reported with the
    fit's error bounds.
  \item \textbf{Structural.} A property of request geometry, independent of any
    provider response.
\end{itemize}

The distinction is load-bearing for the results below: Section 3 is measured,
Section 4 is estimated, and a reader who conflates them will over-read
Section 4.

\section{A Taxonomy of Cache Failures}

Over $C_{\text{cause}}$ the engine isolates 735 breaks accounting for
31{,}264{,}349 re-billed tokens. Table~\ref{tab:breakdown} gives the
distribution. Figures in this section are \emph{measured}.

\begin{table}[h]
\centering
\caption{Re-billed tokens by cause, $C_{\text{cause}}$ (735 breaks,
31.26M tokens).}
\label{tab:breakdown}
\vskip 0.1in
\small
\begin{tabular}{lrrr}
\toprule
\textbf{Cause} & \textbf{Breaks} & \textbf{Share} & \textbf{Mean} \\
\midrule
Client re-render   & 583 & 50.8\% & 27{,}232 \\
TTL expiry         & 39  & 33.9\% & 271{,}436 \\
Prefix divergence  & 102 & 7.0\%  & 21{,}503 \\
Definition change  & 5   & 5.8\%  & 361{,}400 \\
Model switch       & 6   & 2.6\%  & 133{,}667 \\
\bottomrule
\end{tabular}
\end{table}

The distribution's shape matters more than its ranking. Client re-renders are
frequent and individually small; TTL expiries are rare and individually
enormous, at roughly ten times the mean size. One developer leaving for lunch
costs more than a hundred re-renders. Grouped by the intervention that could
plausibly address them, prefix layout and ordering account for 63.6\% of
re-billed tokens and timing for 33.9\%.

\subsection{A sampling failure worth reporting}

This measurement was first run over the 40 largest sessions in the same corpus
and produced approximately the opposite result: TTL expiry at 75.2\% against
33.9\%, and client re-render at 2.5\% against 50.8\%. The conclusion drawn from
the sample --- that three quarters of the recoverable spend is idleness, and
that prefix layout is not worth building --- reverses on the full corpus.

The bias is mechanical. The largest sessions are the longest-running; the
longest-running contain the longest gaps; and a long gap is what TTL expiry
\emph{is}. Sorting by size selected for the cause. We report this because the
sampled figure was defensible on its face, and nothing internal to it would
have exposed the error.

\section{Falsifying Keepalive for This Workload}

A common mitigation is to issue periodic requests that hold the cache open.
Prior work derives a break-even horizon for the practice,
$I_{\max} \approx \tau(w/r - 1)$, with $\tau$ the ping interval, $w$ the
re-prefill multiple and $r$ the cached-read multiple, and reports a paying band
of roughly 46 minutes at $r=0.10$.

Two observations follow, both on $C_{\text{gap}}$. Figures in this section are
\emph{estimated}.

\subsection{The band has moved}

The published band assumes $r = 0.10$. Current top-tier models in our compiled
price table carry $r = 0.025$, which widens $I_{\max}$ to approximately 196
minutes at the same $\tau$. Nothing in the derivation is wrong; a constant moved
underneath it. This is the ordinary failure mode of a published constant, and
it argues for tools that read a live table rather than embed one.

\subsection{The band is nearly empty}

Bucketing all 26{,}675 consecutive assistant-turn pairs by the gap preceding
them, and pricing the counterfactual saving as
$\text{creation} \times \text{input} \times (w - r)$:

\begin{table}[h]
\centering
\caption{Cache-creation spend by preceding gap, $C_{\text{gap}}$.}
\label{tab:gaps}
\vskip 0.1in
\small
\begin{tabular}{lrrr}
\toprule
\textbf{Gap band} & \textbf{$n$} & \textbf{Creation} & \textbf{Recoverable} \\
\midrule
$<$ 5 min            & 26{,}022 & 98.5M & none \\
5--46 min            & 598 & 2.1M & \$11.65 \\
46--196 min          & 40 & 11.4M & \$65.59 \\
$>$ 196 min          & 15 & 8.7M & none \\
\bottomrule
\end{tabular}
\end{table}

Recoverable by keepalive: \$77.24. Not recoverable: \$508.83.

87\% of cache-creation spend on this corpus occurs after gaps shorter than five
minutes, during which no published TTL had elapsed. The cache was alive and was
invalidated anyway --- something rewrote the prefix while it was still warm. No
ping interval reaches that spend, because there is no idle period to bridge.
Keepalive is a real mechanism answering a question this workload does not ask,
and it is the wrong lever here by roughly $7\times$.

We state the scope plainly: this is one workload class on one machine, at the
estimated tier. It falsifies the keepalive hypothesis \emph{for this corpus}. It
does not establish that keepalive fails for workloads with genuine idle
structure, and the cited derivation remains correct on its own terms.

\section{Fan-Out and the Isolation Requirement}

Agentic systems fan out. Sub-agents issue requests concurrently against a shared
session identifier, and their prefixes diverge from one another by design.

An early build of the engine compared each request's prefix hash against a
single session-wide field. Under fan-out this compares a request against
whichever sibling wrote last, so every lane appears to have had its prefix
rewritten by every other lane. The instrument reported that 98.8\% of re-billed
tokens in a 30-lane trial were attributable to changed system or tool
definitions. The figure reached a commit message and a source comment as
established fact before it was identified as an artifact of the comparison
rather than a property of the workload.

The correction is per-lane isolation: a prefix may only be compared against its
own lane's predecessor. On a corrected six-minute run over $C_{\text{lane}}$,
336{,}060 of 7{,}996{,}448 billed prompt tokens --- 4.2\% --- were re-billed by
cache break.

\textbf{These two percentages do not share a denominator and must not be read
as a before-and-after of one quantity.} The 98.8\% figure is a share of the
uncorrected run's own deficit total, itself a product of the faulty comparison.
The 4.2\% figure is re-billed tokens against every prompt token billed. The
first number is not a worse estimate of the second; it is a different quantity,
and reporting the pair as a correction would misstate both. We note this
explicitly because a draft of this paper did exactly that.

\section{Provenance as an Engineering Constraint}

The preceding sections each depend on a discipline that is easier to state than
to hold.

\subsection{Three values, not two}

Absence, zero, and unknown are distinct. A request for which reuse was not
observed is not a request that reused nothing, and collapsing the two produces a
figure that looks measured and is not. We encountered this concretely in a
local-inference surface whose logs record a reuse counter only on requests whose
prompt was already fully resident. A hit rate computed over those requests
measures a population selected for having been cached, and is bounded near
$(n-1)/n$ by the server's own re-evaluation rule. The engine therefore reports
two counts and no ratio for that surface.

\subsection{A check that cannot fail is not evidence}

We verify assertions by mutation: a guard is disabled, and a test suite that
stays green has not been testing it. The discipline has a failure mode we
record here because it is easy to miss. A mutant the compiler rejects --- for
instance, one that leaves a variable unused --- produces a build failure, not a
test failure. An automated harness that greps for test-failure output counts
that build failure as a pass and reports the mutant as caught. We have
mis-scored our own guards this way. A compiler-rejected mutant is not evidence
about a test; it is evidence that the mutation was ill-formed.

\section{Threats to Validity}

\textbf{Single operator, single machine.} All four corpora derive from one
engineer's work. The cause distribution in Section 3 is a property of this
workload and this client, not of agentic coding generally. A client that
re-renders history differently would move the largest row in
Table~\ref{tab:breakdown}.

\textbf{Tier mixing.} Section 3 is measured; Section 4 is estimated from a
byte-to-token fit. The 87\% figure inherits that fit's error, and we do not
claim measured precision for it.

\textbf{Unsolicited transcripts.} The corpora were not designed. This removes
workload-shaping bias and introduces coverage gaps: causes absent from this
work are absent from the taxonomy, which is why \emph{unknown} is a reported
category rather than a residual folded into the nearest named cause.

\textbf{Seat freshness of the price table.} Section 4's conclusion depends on
$r$, which providers change. The 196-minute band is correct for the table
compiled at the date given and for no other.

\section{Conclusion}

Cost claims about agentic systems require turn-level differential attribution.
Aggregate consumption tracking cannot distinguish a session that is large from
a session that is re-billing, and the two call for opposite interventions.

On the workload measured here, the recoverable spend is structural rather than
temporal: prefix layout accounts for 63.6\% of re-billed tokens, and 87\% of
cache-creation spend occurs inside windows where no cache had expired. The
engineering implication is that context layout, not cache timeout tuning, is
where the money is.

We report two of our own measurement failures alongside the results --- a
sampling bias that inverted a conclusion, and an instrument artifact that
produced a confident and meaningless 98.8\% --- because both were internally
consistent and neither would have been caught by a reader who trusted the
output. The tooling around a measurement is part of the measurement.

\section*{Availability}

Source, dated evidence files for every figure above, and the architectural
decision records referenced in Section 6 are published at
\url{https://github.com/RedRobotKK/Replay} under the Business Source License
1.1, converting to Apache 2.0 on 2029-09-06. The engine runs locally, requires
no account, and issues no network request that the operator did not invoke.

\small
\begin{thebibliography}{9}
\bibitem{keepalive} \emph{Keeping the Cache Warm Pays: Keepalive Economics for
  Agentic Workloads.} arXiv:2607.19214. Derives the break-even horizon
  $I_{\max} \approx \tau(w/r - 1)$ evaluated in Section 4.
\bibitem{dontbreak} \emph{Don't Break the Cache: An Evaluation of Prompt
  Caching for Long-Horizon Agentic Tasks.} arXiv:2601.06007v2. Reports cost
  savings from excluding dynamic tool results from the cached prefix.
\bibitem{tracelab} arXiv:2606.30560. Approximately 4{,}300 real Claude Code and
  Codex sessions, approximately 350k steps, publicly released; reports high but
  imperfect prefix cache hit rates.
\bibitem{tracecopilot} arXiv:2608.00101. Production traces across 3.2M users:
  cache hit rate 90\% within a turn and 55\% across turn boundaries, reported as
  drastically invalidated by model switches and compaction.
\end{thebibliography}

\end{document}
